\documentclass[11pt]{article}
\usepackage[margin=25mm]{geometry}
\usepackage{fontspec}
\setmainfont{Latin Modern Roman}
\usepackage{amsmath,amssymb,mathtools}
\usepackage{unicode-math}
\setmathfont{Latin Modern Math}
\usepackage{xurl,hyperref}
\hypersetup{colorlinks=true,urlcolor=blue}
\setcounter{secnumdepth}{0}
\setlength{\parindent}{0pt}
\setlength{\parskip}{5pt}
\setlength{\emergencystretch}{3em}
\newcommand{\sref}[2]{\hyperlink{source:#1:#2}{[S#2]}}
\newcommand{\cataloguescope}[1]{\paragraph{Recorded target.}#1}
\newcommand{\R}{\mathbb R}
\newcommand{\C}{\mathbb C}
\newcommand{\N}{\mathbb N}
\newcommand{\Z}{\mathbb Z}
\newcommand{\Q}{\mathbb Q}
\newcommand{\F}{\mathbb F}
\newcommand{\eps}{\varepsilon}
\begin{document}
\section{$u$-Equations for Arbitrary Hyperbolic Surfaces}\label{statement}

\subsection{Definitions and mathematical statement}

\paragraph{Shared notation.}$\mathbb N=\{1,2,\ldots\}$; the standard number fields have their usual meanings. All additional conventions are stated locally.


A hyperbolic surface has constant negative-curvature geometry; boundary components and marked cusps are retained as part of its type. The report's $u$-equations are algebraic relations among coordinates attached to compatible surface or cluster data. Extension means defining analogous coordinates and relations beyond the surface classes already treated.

\paragraph{Catalogue wording retained for traceability.}
Do the proposed $u$-equations extend to arbitrary hyperbolic surfaces with boundary and marked cusps?

\paragraph{Scope and source qualification.}The $u$-equations refer to the source construction, not an unspecified scalar differential equation. The complete coordinate presentation and admissible surface conventions must accompany an extension.
\subsection{Short English statement}
Determine whether the proposed coordinate equations extend to every hyperbolic surface with boundary and marked cusps.
\subsection{Research attempt}
\paragraph{Model, date and investigated scope.}
OpenAI GPT-6 Astra (configured model \texttt{gpt-6-astra}), ultra reasoning; generated 11 October 2026. Proved infinitely many embedded winding-arc classes on an annulus, exposing the failure of a finite all-arc enumeration shortcut.

\paragraph{Sources and literature check.}
Read the supplied definition, catalogue wording, source qualifications and reference entries for this question. The historical source/context and inherited status are retained; no fresh literature or novelty certification is claimed. No external theorem is used without proof in this attempt; explicit model assumptions are distinguished from facts about the source model.

\paragraph{An infinite-arc obstruction to a finite-type shortcut.} Consider an annulus with one marked point on each boundary. Its universal cover is a strip \(\mathbb R\times[0,1]\), with deck translation \((x,y)\mapsto(x+1,y)\). For each integer \(n\), project the straight segment joining \((0,0)\) to \((n,1)\). It is an embedded arc joining the marked points: its interior has strictly increasing second coordinate, so two distinct interior parameters cannot differ by a deck translation.
These arcs are distinct up to endpoint-fixed homotopy. A homotopy lifts with initial endpoint \((0,0)\); the terminal endpoint must stay in the discrete set \(\mathbb Z\times\{1\}\), hence its integer label is constant. Thus
\[
 \{\text{relative arc classes}\}\supseteq\{\gamma_n:n\in\mathbb Z\}.
\]
\paragraph{Consequence for an attempted extension.} A construction that assigns one variable to every arc cannot be extended to this boundary type by simply retaining a finite catalogue of all arc labels. A finite seed might still generate infinitely many variables, or a completion may handle infinite relations, but that is additional structure. Counting a finite triangulation alone does not remove the winding classes.
\paragraph{Source comparison.} The actual source \(u\)-equations and compatibility exponents are not replaced by an invented differential equation or arbitrary product formula. The first gap is specifying their coordinate presentation and mutation/gluing consistency on surfaces with infinitely many relevant arc classes, then proving the desired relations. The annulus topology is a diagnostic for arbitrary-boundary surfaces, not a counterexample to every possible extension or a new family of verified \(u\)-equations.

\paragraph{Verification and adversarial audit.}
Checked embedded projected arcs and invariance of their lifted terminal integer; distinguished finite seeds from a finite total arc set. This is a GPT-6 Astra Ultra self-audit, not independent human verification. Prompt A drafting followed by prompt B scope/dependency/boundary self-audit. Checked embedded projected arcs and invariance of their lifted terminal integer; distinguished finite seeds from a finite total arc set. The remaining gap is: Extend the actual source coordinate and compatibility relations consistently to arbitrary surface types with infinite arc classes. Unresolved status and zero original-target completion are retained.

Coordinator checked embeddedness of the lifted winding arcs and invariance of the lifted endpoint under relative homotopy. This additional check was model-assisted, not independent human verification.

\paragraph{Tools, process and contributions.}
Python 3 and local file-reading/writing tools for manuscript serialization and structural checks; mathematical calculations were symbolic reasoning, with no numerical experiment or theorem prover. Contribution prompt A was followed by prompt B. GPT-6 Astra with ultra reasoning drafted the research and performed its stated self-audit; any additional coordinator contribution is identified in the verification and audit. A GPT-6-based Codex coordinator assembled the manuscript and checked repository integration; its exact variant and reasoning configuration were not exposed. No human mathematical contribution, independent human proof check or formal proof certification is supplied. Complete derivations and the consulted source references are included here; private reasoning transcripts are not published.

\paragraph{Final status and first remaining gap.}
\textbf{UNRESOLVED}. Extend the actual source coordinate and compatibility relations consistently to arbitrary surface types with infinite arc classes.

\section{Completion Estimate}
\noindent\textbf{COMPLETION: 0\%}
The proved calculation is an elementary auxiliary result or a test of a restricted strategy; no new portion of the full source target is established.

\subsection{Sources}
\begin{itemize}
\item[\hypertarget{source:statement:1}{S1}] ULAM.AI and the UnsolvedMath Contributors, UnsolvedMath, catalogue entry. Catalogue curation under CC BY 4.0; definitions and English exposition added locally.
\url{https://huggingface.co/datasets/ulamai/UnsolvedMath}.
\item[\hypertarget{source:statement:2}{S2}] Mini-Workshop: Scattering Amplitudes, Cluster Algebras, and Positive Geometries (2021), Oberwolfach Reports, DOI 10.4171/owr/2021/57, pp. 3101--3133. Original report and mathematical context.
\url{https://doi.org/10.4171/owr/2021/57}.
\end{itemize}
\label{end:statement}
\end{document}
